\section{Diseño de la Interfaz (Frontend Design)}
Antes de comenzar con la implementación técnica, se elaboró un diseño visual exhaustivo centrado en la usabilidad y la reducción de la fatiga visual. Para ello, se optó por un esquema de colores en modo oscuro (\textit{Dark Mode}), el cual es el estándar en entornos de desarrollo modernos (como VS Code o IntelliJ).

El diseño prioriza un área de chat amplia, donde los mensajes del tutor socrático y del alumno se distinguen claramente por colores y alineación. Además, incluye un panel lateral colapsable para acceder al historial de sesiones anteriores. Se crearon diversos \textit{mockups} y \textit{wireframes} en Figma para validar estas decisiones de UX/UI antes de escribir el código fuente.

\section{Desarrollo del Frontend}
El desarrollo técnico del proyecto comenzó por el Frontend. Esta decisión se tomó para poder interactuar lo antes posible con una interfaz tangible y realizar pruebas de usabilidad tempranas con usuarios reales (Mocking the AI responses). Una vez validada la experiencia de usuario y la correcta visualización de los fragmentos de código renderizados, se procedió a construir el motor de inteligencia artificial que soportaría dicha interfaz.

El frontend se ha implementado utilizando la biblioteca React.js junto con tailwindcss, garantizando una alta responsividad en cualquier tipo de dispositivo y un desarrollo ágil basado en componentes reutilizables.

\section{Desarrollo del Backend y de la API}
Habiendo consolidado la interfaz visual, el backend se programó en Node.js, implementado sobre el framework Express. Para la lógica del tutor socrático, se integró el SDK oficial de OpenAI. Al inicializar un hilo (\textit{thread}) de conversación nuevo en la base de datos PostgreSQL, se envía un \textit{system prompt} invisible para el usuario que condiciona toda la interacción futura del modelo.

\section{Ejemplo de Diálogo}
Se expone un caso de uso real extraído de las pruebas funcionales:
\begin{itemize}
    \item \textbf{Usuario:} ``¿Por qué me da error esta función en Python? \texttt{def suma(a, b) return a+b}''
    \item \textbf{Tutor Socrático:} ``Analiza la primera línea. Tras los argumentos `a` y `b`, ¿falta algún signo de puntuación en Python para iniciar el bloque del código de la función?''
\end{itemize}
Como se detalla en este caso, se cumple con la misión docente de no facilitar el código para que sea el propio alumno quien lo descubra y lo corrija empíricamente.
